


\documentclass[11pt]{article}
\usepackage[letterpaper,margin=1in]{geometry}
\usepackage{amsmath, amssymb, amsthm}
\usepackage{enumitem}
\usepackage{fancyhdr}

\setlength{\headheight}{14pt}
\pagestyle{fancy}
\fancyhf{}
\fancyhead[L]{\textbf{Notes \& Practice}}
\fancyhead[R]{Algebra 1 --- Transformations}
\fancyfoot[C]{\thepage}

\newcommand{\sol}{\par\noindent\textit{Solution.}\;}

\title{\vspace{-0.5in}\textbf{Transformations of Functions}\\[2pt]
       \large Notes with Similar Practice Problems}
\date{}
\author{}

\begin{document}
\maketitle
\vspace{-0.4in}


\section*{Transformations}

A \textbf{transformation} changes the position, shape, or orientation of a graph.

There are several major types of transformations:

\begin{itemize}

\item Horizontal Shifts
\item Vertical Shifts
\item Horizontal Stretches and Shrinks
\item Vertical Stretches and Shrinks
\item Reflections

\end{itemize}

\subsection*{1. Horizontal Shifts}

\[
y=f(x-h)
\]

moves the graph:

\[
\text{right }h\text{ units}
\]

\[
y=f(x+h)
\]

moves the graph:

\[
\text{left }h\text{ units}
\]

\subsection*{2. Vertical Shifts}

\[
y=f(x)+k
\]

moves the graph upward by:

\[
k\text{ units}
\]

\[
y=f(x)-k
\]

moves the graph downward by:

\[
k\text{ units}
\]

\subsection*{3. Horizontal Stretch and Shrink}

\[
y=f(ax)
\]

causes a horizontal shrink by a factor of:

\[
\frac1a
\]

The $x$-coordinates divide by $a$.

If:

\[
(x,y)
\]

is on the graph of:

\[
y=f(x)
\]

then:

\[
\left(\frac{x}{a},y\right)
\]

is on the graph of:

\[
y=f(ax)
\]

\vspace{0.3cm}

Similarly,

\[
y=f\left(\frac{x}{a}\right)
\]

causes a horizontal stretch by a factor of:

\[
a
\]

\subsection*{4. Vertical Stretch and Shrink}

\[
y=af(x)
\]

causes a vertical stretch by a factor of:

\[
a
\]

The $y$-coordinates multiply by $a$.

\vspace{0.3cm}

\[
y=\frac1a f(x)
\]

causes a vertical shrink by a factor of:

\[
\frac1a
\]

\subsection*{5. Reflections}

\[
y=-f(x)
\]

reflects the graph across the $x$-axis.

\vspace{0.3cm}

\[
y=f(-x)
\]

reflects the graph across the $y$-axis.

\section*{Summary}

\subsection*{Most Important Ideas}

\[
\boxed{
\text{Inside changes affect }x\text{-values inversely}
}
\]

\[
\boxed{
\text{Outside changes affect }y\text{-values directly}
}
\]

\subsection*{Horizontal Transformations}

\[
f(ax)
\]

means:

\[
x\rightarrow\frac{x}{a}
\]

\begin{itemize}

\item $f(2x)$: horizontal shrink by $\frac12$
\item $f(3x)$: horizontal shrink by $\frac13$
\item $f\left(\frac{x}{2}\right)$: horizontal stretch by $2$
\item $f\left(\frac{x}{5}\right)$: horizontal stretch by $5$

\end{itemize}

\subsection*{Vertical Transformations}

\[
af(x)
\]

means:

\[
y\rightarrow ay
\]

\begin{itemize}

\item $2f(x)$: vertical stretch by $2$
\item $5f(x)$: vertical stretch by $5$
\item $\frac12 f(x)$: vertical shrink by $\frac12$

\end{itemize}

\subsection*{Coordinate Mapping Rules}

\[
y=f(ax)
\]

maps:

\[
(x,y)\rightarrow\left(\frac{x}{a},y\right)
\]

\vspace{0.3cm}

\[
y=af(x)
\]

maps:

\[
(x,y)\rightarrow(x,ay)
\]

\subsection*{Key Principle}

Any point satisfying the equation is on the graph, and every point on the graph satisfies the equation.

\newpage

\section*{Practice Problems}

\begin{enumerate}

\item Describe the transformation:

\[
y=f(x-4)
\]

\item Describe the transformation:

\[
y=f(x+7)
\]

\item Describe the transformation:

\[
y=f(x)+5
\]

\item Describe the transformation:

\[
y=f(x)-3
\]

\item Describe the transformation:

\[
y=f(2x)
\]

\item Describe the transformation:

\[
y=f(5x)
\]

\item Describe the transformation:

\[
y=f\left(\frac{x}{2}\right)
\]

\item Describe the transformation:

\[
y=f\left(\frac{x}{6}\right)
\]

\item Describe the transformation:

\[
y=3f(x)
\]

\item Describe the transformation:

\[
y=\frac12 f(x)
\]

\item Describe the transformation:

\[
y=-f(x)
\]

\item Describe the transformation:

\[
y=f(-x)
\]

\item If $(8,2)$ is on the graph of $y=f(x)$, find the corresponding point on:

\[
y=f(4x)
\]

\item If $(3,9)$ is on the graph of $y=f(x)$, find the corresponding point on:

\[
y=f\left(\frac{x}{3}\right)
\]

\item If $(6,-1)$ is on the graph of $y=f(x)$, find the corresponding point on:

\[
y=2f(x)
\]

\item If $(10,5)$ is on the graph of $y=f(x)$, find the corresponding point on:

\[
y=\frac12 f(x)
\]

\item If $(4,7)$ is on the graph of $y=f(x)$, find the corresponding point on:

\[
y=-f(x)
\]

\item If $(-3,2)$ is on the graph of $y=f(x)$, find the corresponding point on:

\[
y=f(-x)
\]

\item Given the parent function:

\[
y=x^2
\]

graph:

\[
y=f(2x)
\]

and identify whether the graph is stretched or shrunk.

\item Explain why:

\[
f(3x)
\]

causes a horizontal shrink by a factor of:

\[
\frac13
\]

using coordinate mapping.

\end{enumerate}
\end{document}